% ===================================================================
%  TÁBLA Uimh. 10 — An Fharraige. — An Calafort.
%  Traduction 2026 d'apres texto/io, controlee sur texto/fr, texto/en
%  et texto/es. Consigne : voir texto/ga/10-tabla-01.tex.
%
%  Le francais decale sa numerotation d'un cran a partir de l'alinea 8
%  — coquille de l'imprimeur francais. L'irlandais suit l'ido.
%
%  LE TABLEAU DE LA MER EST LE PLUS EMPRUNTE DE LA COLONNE, comme il
%  l'a ete en ukrainien, en basque et en roumain, et pour la meme
%  raison : la marine a vapeur est arrivee toute faite au XIXe siecle,
%  avec ses mots. « frigéad », « toirpéad », « cúrsóir », « píolóta »,
%  « lián », « canáil », « seamafór », « aimiréal » — l'irlandais les
%  prend tels quels et les decline. Ce qui reste celtique, ce sont les
%  mots que l'on apprend DEHORS et sans maitre : « trá » la plage,
%  « aill » la falaise, « cladach » le rivage, « tonn » la vague,
%  « caolas » le detroit, « oileán » l'ile, « crann » le mat, « rámha »
%  l'aviron, « líonta » les filets, « cé » le quai. Le partage est le
%  meme dans les quatre langues, et c'est le lieu de l'apprentissage
%  qui le trace : la mer que l'on regarde garde ses mots, la marine
%  que l'on construit prend ceux qui viennent avec les navires.
%
%  UN MOT QUI VAUT POUR DEUX. Le tableau nomme deux fois l'ouvrage ou
%  l'on accoste : la « warfo » (14), qui est la jetee, et le « kayo »
%  (58), qui est le quai. L'irlandais courant dirait « cé » des deux
%  cotes, et le lecteur perdrait la difference que la planche grave.
%  On ecrit donc « lamairne » pour la jetee et « cé » pour le quai.
%
%  DEUX RENVOIS QUE LE VSO AURAIT DEPLACES si l'on avait suivi l'ordre
%  ordinaire de la phrase : « potenta krani (62) deprenis li » et
%  « esas fuorto (97) tre importanta ». Le verbe irlandais precede son
%  sujet, et le sujet porte ici le renvoi ; l'objet, lui, n'en porte
%  pas. L'inversion ne se produit donc pas — elle exige DEUX elements
%  qui portent chacun un renvoi. Le tableau 10 le montre par l'absence
%  la ou les tableaux 7 et 9 l'ont montre par la presence.
% ===================================================================

%%K t10-apar-1 apar
\VUsaut{4.0mm}
\VUtitre{15.0pt}{60}{TÁBLA Uimh. 10}
\VUsaut{3.0mm}
\VUpk{11.4pt}{40}{An Fharraige. --- An Calafort.}
\VUsaut{3.0mm}

%%K t10-01-1 p
1. --- Sa samhradh, agus cuid againn ag lorg an chiúnais agus na fionnuaire ar an
sliabh, is fearr le daoine eile dul go dtí an cladach. Sin é a rinneamar anuraidh,
mar is mór againn \VUgras{an tAigéan} \textsuperscript{(1)}.

%%K t10-02-1 p
2. --- Maidir liom féin, is aoibhinn liom breathnú ar an leathanas ollmhór uisce sin
a shíneann amach go \VUgras{bun na spéire} \textsuperscript{(2)}, agus na
\VUgras{galtáin} ollmhóra \textsuperscript{(3)} a fheiceáil ag imeacht ar
\VUgras{aistear} fada leis na taistealaithe a bhíonn ag dul go Meiriceá.

%%K t10-03-1 p
3. --- Dá bhrí sin roghnaíonn m'athair, a bhfuil eolas aige ar mo dhúil, trá
an-chiúin i gcónaí chun na laethanta saoire a chaitheamh, trá atá suite in aice le
ceann de na \VUgras{calafoirt chogaidh} mhóra atá againn.

%%K t10-03-2 p
Le linn ár dtréimhse dheireanaí tháinig an \VUgras{cabhlach} ar ancaire taobh thiar
den \VUgras{mhúr farraige} ollmhór \textsuperscript{(4)} a dhúnann agus a chosnaíonn
an \VUgras{calafort cogaidh} \textsuperscript{(5)}, agus d'éirigh liom, mar ba mhian
liom, na \VUgras{longa cogaidh} éagsúla a scrúdú: an \VUgras{fomhuireán} beag
\textsuperscript{(10)} a thumann agus a imíonn as radharc faoi na huiscí, agus fós an
\VUgras{chathlong} ollmhór \textsuperscript{(9)} nach raibh eagla uirthi go dtí seo ach
roimh ionsaithe an \VUgras{bháid toirpéad} \textsuperscript{(8)}.

%%K t10-tit-2 sub
\VUsaut{3.0mm}
\VUpk{10.2pt}{40}{Na longa beaga.}
\VUsaut{3.0mm}

%%K t10-04-1 p
4. --- Bhí a fhios ag an gceann sin cheana conas an pointe lag san armúr tiubh
miotail a aimsiú go \VUgras{healaíonta} chun an \VUgras{toirpéad} contúirteach a chur
ann, toirpéad a chuireann an long ar lár lena phléasc; ach ba lú ba chás é ná an
fomhuireán, mar leanann na scriostóirí é gan dua.

%%K t10-05-1 p
5. --- Chonaic mé freisin na \VUgras{cúrsóirí} mearluatha \textsuperscript{(6)}
armúrtha, atá beagnach chomh doghonta leis an bhfíorchathlong, roinnt \VUgras{long
iompair} \textsuperscript{(12)}, trí nó ceithre \VUgras{bhád gunnaí}
\textsuperscript{(7)}, agus \VUgras{frigéad} \textsuperscript{(11)} ar deireadh.
\VUgras{Is é radharc an chabhlaigh sin, agus é ag ullmhú chun na hancairí a thógáil,
an ceann is suimiúla ar fad}.

%%K t10-06-1 p
6. --- A leithéid de dhifríocht tógála idir na carnáin ollmhóra cruach sin agus na
\VUgras{trímhastaigh} ghalánta nó na \VUgras{longa seoil} caola \textsuperscript{(13)}
a bhíonn fós in úsáid sa chabhlach tráchtála, longa a chonaic mé ag teacht isteach sa
chalafort faoi lánseol nuair a bhínn ag siúl ar an \VUgras{lamairne}
\textsuperscript{(14)}. Is fíor gur chaill siad seo cuid mhór dá mbuntáistí
\VUgras{nuair} a bhí an \VUgras{ghaoth} ina n-aghaidh. B'éigean dóibh na seolta go
léir a ísliú chun dul isteach sa duga arís, agus bhí \VUgras{tuga}
\textsuperscript{(16)} riachtanach chun iad a thógáil agus a tharraingt, mar a bheadh
\VUgras{báirsí} simplí \textsuperscript{(17)} ann, chun na cé.

%%K t10-tit-3 sub
\VUsaut{3.0mm}
\VUpk{10.2pt}{40}{An calafort tráchtála.}
\VUsaut{3.0mm}

%%K t10-07-1 p
7. --- Is é an rud a chuireann suim sa chalafort a bhfuil mé ag caint air, nach
calafort cogaidh amháin atá ann, a rinneadh de láimh le hoibreacha ollmhóra, ach
\VUgras{calafort tráchtála} iontach chomh maith, atá suite in \VUgras{inbhear}
abhann móire.

%%K t10-08-1 p
8. --- Tá an teanga talún a scarann an dá dhuga daingnithe agus cosanta ag sraith
\VUgras{ceallraí} agus \VUgras{báisteán} a bhrúnn amach ar an bhfarraige. Teastaíonn
cead ar leith chun siúl ar na \VUgras{rampair} \textsuperscript{(15)}. Fuair mé é sin
gan dua trí m'uncail, atá ina innealtóir ar na \VUgras{longchlóis}
\textsuperscript{(18)}, agus chuaigh mé ag féachaint ar na longa a thógtar faoi
sciobóil fhairsinge.

%%K t10-09-1 p
9. --- Bhí mé i láthair fiú ag seoladh cathloinge, agus is cuimhin liom an nóiméad
corraitheach a mhothaigh an slua ar fad nuair a thosaigh an mais ollmhór, fágtha
fúithi féin, ag sleamhnú ar a creat sleamhnáin agus tháinig sí ar snámh ar uisce na
habhann.

%%K t10-10-1 p
10. --- Chun dul chuig na longchlóis trasnaítear an \VUgras{machaire cleachtaí}
\textsuperscript{(19)} ina dtagann saighdiúirí an gharastúin gach lá chun a n-oiliúint
a dhéanamh, agus téitear thar \VUgras{dhroichead coise} iarainn
\textsuperscript{(20)} a cheanglaíonn faiche na paráide leis an \VUgras{armlann}
\textsuperscript{(21)}.

%%K t10-tit-4 sub
\VUsaut{3.0mm}
\VUpk{10.2pt}{40}{An Armlann agus na Dugaí.}
\VUsaut{3.0mm}

%%K t10-11-1 p
11. --- Thug mé cuairt le m'uncail ar an bhfoirgneamh mór sin atá lán de
\VUgras{ghunnaí móra} \textsuperscript{(22)}, de \VUgras{urchair mhóra}
\textsuperscript{(23)} agus de \VUgras{phléascáin}. Chonaic mé freisin an
\VUgras{mhonarcha mhiotail} \textsuperscript{(24)} ina ndéantar \VUgras{coirí}
\textsuperscript{(25)} na ngaltán.

%%K t10-12-1 p
12. --- Go han-ghar dó tá na \VUgras{dugaí} \textsuperscript{(26)} --- dugaí
deisiúcháin --- agus na \VUgras{dugaí ar snámh} \textsuperscript{(27)}
sár-shuimiúla, inar éirigh liom, ar long a bhí á deisiú, an \VUgras{lián}
\textsuperscript{(28)}, an \VUgras{chíle} \textsuperscript{(29)} agus \VUgras{stiúir}
\textsuperscript{(30)} báid a fheiceáil go gairid.

%%K t10-13-1 p
13. --- Tá an calafort tráchtála chomh fiúntach le staidéar leis an gcalafort
cogaidh, agus chaith mé mórán uaireanta ag stánadh ar na céanna ina mbíonn na
\VUgras{báid rothacha} \textsuperscript{(31)} a théann suas an abhainn ar ceangal,
agus ag féachaint ar an \VUgras{gcraein chrannaithe} \textsuperscript{(32)} ag obair,
craein a thógann ualaí ollmhóra mar a bheadh cleití ann.

%%K t10-tit-5 sub
\VUsaut{3.0mm}
\VUpk{10.2pt}{40}{An Píolóta.}
\VUsaut{3.0mm}

%%K t10-14-1 p
14. --- Bhí meas agam ar na \VUgras{longa trasatlantacha} \textsuperscript{(34)} ag
imeacht as an ród, agus a \VUgras{mbratacha} \textsuperscript{(35)} in airde san aer,
faoi stiúir \VUgras{phíolóta} oilte \textsuperscript{(36)} a bhí in ann an
\VUgras{cainéal} a leanúint go hiontach, cainéal atá marcáilte le \VUgras{baoithe}
\textsuperscript{(40)} agus le \VUgras{comharthaí}, agus é ag seachaint na
\VUgras{ngeabaraí} \textsuperscript{(41)} a stiúrann iad féin go doiligh lena n-aon
seol cearnach amháin. D'aithin mé ar an \VUgras{gcuid tosaigh}
\textsuperscript{(37)} --- an gob --- na \VUgras{cábáin} \textsuperscript{(39)} den
dara grád, agus ar an \VUgras{gcuid deiridh} \textsuperscript{(38)} --- an deireadh
--- na hallaí do phaisinéirí an chéad ghráid.

%%K t10-15-1 p
15. --- Uaireanta bhí mé i láthair freisin ag luchtú \VUgras{loinge lasta}
\textsuperscript{(42)} inar carnadh earraí an \VUgras{stórais} \textsuperscript{(43)}
agus an \VUgras{stáisiúin mhara} \textsuperscript{(44)}, earraí a thug traenacha
líonmhara \VUgras{líne na gcéanna} \textsuperscript{(45)} isteach.

%%K t10-tit-6 sub
\VUsaut{3.0mm}
\VUpk{10.2pt}{40}{An Rámhaíocht.}
\VUsaut{3.0mm}

%%K t10-16-1 p
16. --- Is minic a chuaigh mé ag bádóireacht sa \VUgras{duga snámha}
\textsuperscript{(53)}, inar tháinig na \VUgras{báid bheaga} \textsuperscript{(46)} ar
foscadh, agus ba mhór an t-áthas dom an \VUgras{rámha} \textsuperscript{(47)} a
láimhseáil. Chuaigh mé suas ar \VUgras{dheic} \textsuperscript{(48)} \VUgras{bháid
seoil} \textsuperscript{(49)}, dhreap mé, le \VUgras{téada} \textsuperscript{(52)}, an
\VUgras{crann} \textsuperscript{(51)} isteach sna \VUgras{slata seoil}
\textsuperscript{(50)}, agus ansin lean mé de mo shiúlóid \VUgras{i yóla}
\textsuperscript{(54)} i measc na \VUgras{mbád iascaigh} troma \textsuperscript{(55)},
áit a dtriomaíonn na \VUgras{líonta} \textsuperscript{(56)}.

%%K t10-17-1 p
17. --- Ar an \VUgras{ché} \textsuperscript{(58)} chuaigh na \VUgras{hearraí}
\textsuperscript{(59)} ba éagsúla thar mo shúile, iad i \VUgras{gcistí}
\textsuperscript{(60)} nó i \VUgras{mburlaí} \textsuperscript{(61)}. Ba iad
\VUgras{lasta} \textsuperscript{(63)} na mbáirsí iad, agus thóg \VUgras{craenacha}
cumhachtacha \textsuperscript{(62)} iad agus leag siad ar \VUgras{thrucailí}
\textsuperscript{(64)} iad, fad is a bhí na \VUgras{tiománaithe lasta} ag déanamh a
ndearbhaithe agus ag íoc na ndleachtanna i \VUgras{dteach an chustaim}
\textsuperscript{(65)}.

%%K t10-18-1 p
18. --- Faoi dheireadh, nuair a shroich mé an \VUgras{droichead casta}
\textsuperscript{(66)} nó an \VUgras{loc} \textsuperscript{(69)}, agus mé ag gabháil
thar an \VUgras{dreidire} \textsuperscript{(68)} a ghlanann an \VUgras{chanáil}
\textsuperscript{(67)}, chuaigh mé i dtír ar an lamairne in aice leis an
\VUgras{seamafór} \textsuperscript{(70)}.

%%K t10-19-1 p
19. --- Is ar an taobh eile den lamairne sin a thagann na \VUgras{báid fhada}
\textsuperscript{(71)} i dtír, báid a thugann oifigigh an chabhlaigh chun talaimh. Is
radharc iontach é na mairnéalaigh a fheiceáil ag ardú a rámhaí in éineacht, agus an
bád, á iompar ag an \VUgras{tonn} \textsuperscript{(72)}, ag teacht le taobh dhréimire
an chéibh gan stró. Is san áit seo is fearr a fheictear an \VUgras{taoide} agus
gluaiseacht an \VUgras{láin mhara} agus an \VUgras{laig trá} ar an \VUgras{trá}
\textsuperscript{(73)}.

%%K t10-tit-7 sub
\VUsaut{3.0mm}
\VUpk{10.2pt}{40}{An Cumann.}
\VUsaut{3.0mm}

%%K t10-20-1 p
20. --- Chun filleadh abhaile bhain mé úsáid as an \VUgras{tram leictreach}
\textsuperscript{(74)} a bhfuil a \VUgras{stad} \textsuperscript{(75)} in aice leis an
\VUgras{gcuaille} atá curtha os comhair chumann na n-oifigeach.

%%K t10-20-2 p
Ó \VUgras{bhalcóin} \textsuperscript{(76)} an chumainn, atá maisithe le
\VUgras{hancairí} \textsuperscript{(77)}, le gunnaí móra agus le hurchair, is minic a
chonaic mé cruinniú taibhseach d'oifigigh chabhlaigh: \VUgras{leifteanaint loinge}
\textsuperscript{(78)} agus a gclaíomh leo, \VUgras{captaein airtléire}
\textsuperscript{(79)} agus \VUgras{coisithe} \textsuperscript{(80)} mara agus a
\VUgras{sabar} leo. Taobh thiar díobh bhí \VUgras{mairnéalach}
\textsuperscript{(81)} a thug \VUgras{gloine fhada} chucu nuair ab áil leo an rud a
bhí ag titim amach i gcéin a dhéanamh amach.

%%K t10-21-1 p
21. --- Chonaic mé freisin \VUgras{meirgirí} óga \textsuperscript{(82)} ag glacadh
orduithe óna \VUgras{gceannasaí loinge} \textsuperscript{(83)} nó ón
\VUgras{aimiréal} \textsuperscript{(84)}, fad is a bhí \VUgras{ceathrúshaoiste}
\textsuperscript{(85)} ag fanacht chun iad a bhreith ar bord.

%%K t10-22-1 p
22. --- Ón mbalcóin sin tá an radharc taibhseach: tá an trá ar fad faoi shúil, agus a
\VUgras{pubaill} \textsuperscript{(86)} agus a \VUgras{bothanna} \textsuperscript{(87)}
léi, agus feictear, ar uair an fholctha, na \VUgras{folcadóirí}
\textsuperscript{(88)} ag rince go háthasach os comhair an \VUgras{chasaíno}
\textsuperscript{(89)}.

%%K t10-23-1 p
23. --- Ag ceann na trá déanann an cladach \VUgras{leathinis} bheag
\textsuperscript{(91)} \VUgras{charraigeach}, ar a mbriseann an \VUgras{tonn mhór}
\textsuperscript{(92)}, agus ar ar tógadh \VUgras{teach solais} \textsuperscript{(93)}
a thaispeánann an bealach do na mairnéalaigh san oíche. Níl an leathinis sin
ceangailte leis an talamh ach le \VUgras{cuing talún} an-chúng \textsuperscript{(90)}.

%%K t10-tit-8 sub
\VUsaut{3.0mm}
\VUpk{10.2pt}{40}{Radharc ginearálta ar an gcósta.}
\VUsaut{3.0mm}

%%K t10-24-1 p
24. --- Síneann an \VUgras{aill} \textsuperscript{(94)} léi, í stróicthe ag an
bhfarraige, a tharraingíonn inti go haerach sraith \VUgras{bánna}
\textsuperscript{(95)} agus \VUgras{rinní} (ceanna), rud a fhágann an cladach ar an
gceann is pictiúrtha.

%%K t10-24-2 p
Ar an \VUgras{ardchlár} \textsuperscript{(96)}, a bhfuil radharc uaidh ar an
\VUgras{gcaolas} \textsuperscript{(98)}, tá \VUgras{dún} \textsuperscript{(97)}
an-tábhachtach agus roinnt « villa »-í.

%%K t10-25-1 p
25. --- Scarann an caolas sin ón mórthír \VUgras{oileán} \textsuperscript{(99)}
an-chontúirteach, atá timpeallaithe ag \VUgras{sceireacha} \textsuperscript{(100)}
agus \VUgras{bristeacha}, mar a dtéann báid agus fiú longa móra i dtalamh uaireanta.
In ainneoin luais an fhóirithinte agus theacht mear na \VUgras{mbád tarrthála}, is
uafásach na \VUgras{longbhristeacha} sin, agus le linn \VUgras{stoirmeacha} an
chónocht chuaigh roinnt long ar lár le foireann agus le lasta.

%%K t10-25-2 p
In ainneoin na comharsanachta sin, is minic na mairnéalaigh ag triall ar
\VUgras{an gcalafort}.
